% ============================================================================
% PAPER 2 (method line) -- M1 campaign manuscript SKELETON.
% Working title (NOT final):
%   "Information-optimal illumination for dead-time-limited single-pixel
%    imaging: certificates, water filling, and the jitter-capped ridge"
% ----------------------------------------------------------------------------
% RELATION TO PAPER 1: paper 1 (paper/main_oe.tex, frozen at tag
%   paper1-freeze-v1) mapped WHEN raising flux helps for a given family of hand
%   patterns. Paper 2 asks (a) WHAT illumination extracts the most information
%   from a dead-time-limited detector under fixed physical resources, and
%   (b) HOW HARD a real (jittered) detector can be driven before the ridge
%   collapses. It is cited throughout as \cite{companion} = "[companion, in
%   submission]".
% ----------------------------------------------------------------------------
% PLACEHOLDER DISCIPLINE (inherited, 2026-07-13 audit): every number/verdict
%   that depends on an outstanding output is \SPH{...} and indexed in
%   PLACEHOLDER_LEDGER.md. The manuscript is "filled" only when that ledger is
%   EMPTY. Three fill classes: R14 (theory ruling, issue "R14 ruling:
%   unification architecture"), M1 (confirmatory campaign outputs, tag
%   m1-freeze), USER (author/repo/funding decisions).
% NUMBER PROVENANCE: frozen concepts and constants are transcribed from the
%   rulings named in each block comment (R10 = ROUND63_GPT_ROUND10_RULING_RAW,
%   R11 = ..._ROUND11_..., R13 = ..._ROUND13_..., M1 spec =
%   docs/ROUND63_METHOD_SPEC_M1.md, jitter table =
%   docs/ROUND63_GPT_ROUND14_QUESTION.md). No campaign outcome is asserted.
% ----------------------------------------------------------------------------
% PORTING NOTE: standard article class so it compiles anywhere; port to the
%   target template (opticajnl.cls or an IEEE/TCI class) at submission prep.
% ============================================================================
\documentclass[10pt,twocolumn]{article}
\usepackage{amsmath,amssymb,bm}
\usepackage{graphicx}
\usepackage[margin=2.2cm]{geometry}
\usepackage{xcolor}
\usepackage{hyperref}
\newcommand{\SPH}[1]{\textcolor{red}{[#1]}}   % unfilled campaign placeholder
\newcommand{\rhobar}{\bar{\rho}}
\newcommand{\Cu}{C_{u}}
\newcommand{\rhoR}{\rho_{R}}                    % production ridge target (R11)
\newcommand{\Jex}{J_{\mathrm{exact}}}          % exact per-slot count information
\newtheorem{theorem}{Theorem}
\newtheorem{proposition}{Proposition}

\title{Information-optimal illumination for dead-time-limited
single-pixel imaging:\\ certificates, water filling, and the jitter-capped ridge}

% Author block: USER DECISION (affiliation + collaborators) -- do not fill.
\author{\SPH{author block --- user decision}}
\date{}

\begin{document}
\maketitle

\begin{abstract}
% Field-first opening, then the companion positioning, then the two method
% questions, then the contribution close. NO "first" (R10 Q6 novelty
% discipline). All outcome numbers are \SPH pending the M1 campaign.
Single-pixel imaging with photon-counting detectors is conventionally
operated at low flux to keep dead-time loss small. A companion study mapped
\emph{when} deliberately raising the flux into the nonlinear dead-time
regime helps a fixed family of illumination patterns, and located a
finite-window count-information ridge at the cube-root load
$\rho^{*}(\nu) = (6\nu)^{1/3} - 2/3 + O(\nu^{-1/3})$
($\nu = T/\tau$ dead-time slots per exposure)~\cite{companion}. Here we ask
two follow-up questions:\ what illumination extracts the most information
from a dead-time-limited detector under fixed physical resources, and how
hard a real, timing-jittered detector can be driven before that ridge
collapses. We cast hand-designed sparse and matched patterns as atoms in a
constrained experiment-design problem whose per-atom information is computed
from the exact finite-window nonparalyzable renewal count law, and optimize a
local D-optimal design under explicit detector-load, dose, per-pixel-exposure,
and peak-irradiance constraints, with a generalized-equivalence /
relaxed Kiefer--Wolfowitz certificate relative to a frozen structured
dictionary. The design separates into two conjugate resource corners: a
photon-budgeted variable-load optimizer (\textsc{OED-DT}) and a
time-budgeted, power-available ridge-tracking policy (\textsc{RIDGE-FIXED}).
On the theory side we derive an exact missing-information identity for hidden
random hold times and the long-window count-information rate
$J_\infty(\rho,c_v) = \rho/[(1+\rho)(1+c_v^{2}\rho^{2})]$, whose optimum is
finite for every nonzero hold-time variance and solves the cubic
$c_v^{2}\rho^{2}(1+2\rho) = 1$, scaling as $2^{-1/3}c_v^{-2/3}$ at small
jitter ($c_v$ the hold-time coefficient of variation); matching this
extensive loss to the deterministic terminal-phase loss predicts the smooth
crossover $\rho^{*}(\nu,c_v)\sim(6\nu)^{1/3}(1+12\nu c_v^{2})^{-1/3}$ ---
zero jitter is the critical manifold on which the ridge grows as $\nu^{1/3}$,
and a Monte-Carlo sweep matches the exact-law peak predictions within
$5$--$6\%$ on the measured grid. We report a preregistered simulation campaign (M1) with a single
gate on the optimized arm, a common charged pre-scan, six arms, and a
mandatory per-cell disclosure ledger; all outcomes are reported as frozen
\SPH{M1: primary/secondary verdicts and cross-arm numbers}.
\end{abstract}

% ===========================================================================
\section{Introduction}
% Para 1 -- field-first.
Single-pixel imaging (SPI) and computational ghost imaging reconstruct a
scene from the responses of a single bucket detector to a sequence of
structured illumination patterns~\cite{duarte2008singlepixel,
shapiro2008computational,edgar2019principles}. In the photon-starved regime
the bucket is a photon-counting device that is blind for a dead time $\tau$
after each registered count, and the conventional response is to operate at
low flux, keeping the per-slot load $\rho = \lambda\tau$ at the few-percent
level and paying in acquisition time.

% Para 2 -- position against the companion; state the two questions.
A companion study reframed high-flux operation as an operating-map
question and answered it for fixed hand patterns~\cite{companion}: dense
multiplexed illumination is multiplex-limited, equal-load sparse
illumination yields large but not uniformly broad gains, and the exact
Fisher information of the integrated nonparalyzable count possesses a
principal ridge at $\rho^{*}(\nu) = (6\nu)^{1/3} - 2/3 + O(\nu^{-1/3})$,
beyond which no count-only estimator recovers information. That work took
the illumination family as given. Two questions remain open. First, a
\emph{design} question:\ under a fixed incident-photon and dose budget,
which illumination extracts the most information from the dead-time channel,
and by how much can a certified optimizer improve on a strong fixed
hand design? Second, a \emph{deployment} question:\ the cube-root ridge is
derived for an idealized detector with deterministic dead time; how far into
the ridge can a real detector with timing jitter actually be driven?

% Para 3 -- the approach.
We answer both by treating each hand-designed sparse or matched pattern as an
atom in a constrained local experiment-design problem. Each atom's
information is the exact chain-rule renewal information of Section~\ref{sec:design},
the design is optimized as a concave program over design measures, and the
continuous optimum is certified against the frozen structured dictionary by a
generalized-equivalence / relaxed Kiefer--Wolfowitz bound. The same
formulation exposes a conjugacy: a photon-budgeted optimizer rations a fixed
incident resource across directions (information water filling), whereas a
time-budgeted policy that has power to spare tracks the ridge at every dwell.
A parallel theory line asks how detector timing jitter modifies the ridge
itself.

% Para 4 -- frozen novelty wording (R10 Q6, verbatim intent; no "first").
We formulate a local approximate-design framework for single-pixel
illumination in which each pattern's information is computed from the exact
finite-window nonparalyzable renewal count law, optimized under explicit
detector-load, dose, per-pixel-exposure, and peak-irradiance constraints and
accompanied by generalized-equivalence and exact-rounding certificates.
Existing photon-counting optimal experiment design, adaptive or learned SPI,
and dead-time optimal-flux studies address neighboring
components~\cite{feng2018adaptive,maxu2021sensing,attia2018goaloriented,
burger2021sequential,ariascastro2013adaptive,jorgensen2026deadtime}; we found
no prior framework combining these elements for integrated-count single-pixel
illumination. The certificate is always stated relative to the frozen
dictionary, the local pre-scan estimate, and the declared resource model.

% Para 5 -- contributions.
This work contributes:\ (i)~a master concave design program over design
measures with the exact renewal kernel, and a generalized-equivalence /
relaxed Kiefer--Wolfowitz certificate against a frozen structured dictionary
(Section~\ref{sec:design}); (ii)~a photon-budget versus time-budget
conjugacy that separates the design-allocation question (\textsc{OED-DT})
from the operating-point question (\textsc{RIDGE-FIXED}), with the associated
information-water-filling reading; (iii)~a random-hold count-information theory (Section~\ref{sec:theory}):\ an
exact missing-information identity for hidden iid hold times, the long-window
rate $J_\infty(\rho,c_v) = \rho/[(1+\rho)(1+c_v^{2}\rho^{2})]$ with its exact
jitter-limited optimum $c_v^{2}\rho^{2}(1+2\rho) = 1$ and small-jitter law
$2^{-1/3}c_v^{-2/3}$, and the matched crossover prediction
$(6\nu)^{1/3}(1+12\nu c_v^{2})^{-1/3}$, organized per the frozen unification
ruling with the master design program as framework and the alignment result
as its conditional proposition; and (iv)~a
preregistered simulation campaign (M1) whose scenes, atoms, load grid,
guards, endpoints, and single gate were frozen in an immutable commit before
any confirmatory reconstruction existed, reported as frozen including any
negative primary (Sections~\ref{sec:campaign}--\ref{sec:results})
\SPH{M1: campaign outcomes}. All code, frozen preregistration documents, and
per-cell evidence bundles are openly released.

% ===========================================================================
\section{The design problem}
\label{sec:design}
% Master concave program (R10 Q1/Q4), certificate (R10 Thm R10.1 + R13 6.1),
% conjugacy (R11 4 + R13 8f), guards as constraints (R10 Q4 + R11 7 + R13).

\subsection{Master concave program over design measures}
Freeze the coarse pre-scan reconstruction at $\hat x$ and parameterize the
locally estimable image perturbation as $x(\theta) = \hat x + B\theta$, with
$B \in \mathbb{R}^{n\times r}$ spanning a declared $r$-dimensional estimable
subspace. For a nonnegative illumination atom $a$ the arrival rate is
$\lambda(a) = \Phi\,a^{\top}\hat x + d$ and the load is $\rho(a) = \tau\lambda(a)$.
Let $\Jex(\rho,\nu)$ be the exact per-slot Fisher information of the integrated
nonparalyzable count for $\log\lambda$ derived in the companion
work~\cite{companion}. The score direction for $\theta$ is
$q(a) = \Phi\,B^{\top}a / \lambda(a)$, so the correct rank-one local
information atom is the chain-rule form
\begin{equation}
  \begin{aligned}
  H(a;\hat x) &= \nu\,\Jex(\rho(a),\nu)\,q(a)q(a)^{\top}\\
  &= \nu\,\Jex(\rho(a),\nu)\,\frac{\Phi^{2}}{\lambda(a)^{2}}\,
    B^{\top}a a^{\top}B .
  \end{aligned}
  \label{eq:atom}
\end{equation}
A raw sensitivity $g \propto \Phi a$ is incomplete unless the factor
$\sqrt{\nu \Jex}/\lambda$ is absorbed into it. For a probability measure $\xi$
over the compact feasible atom set $\mathcal{A}$, define
$V(\xi) = V_0 + \int_{\mathcal{A}} H(a)\,d\xi(a)$ with $V_0 \succeq 0$ the
frozen pre-scan information, and maximize the local D-criterion
\begin{equation}
  \max_{\xi}\; \log\det V(\xi)
  \quad\text{s.t.\ linear resource constraints,}
  \label{eq:master}
\end{equation}
subject to average load, cumulative dose, family allocation, and peak
constraints. The objective is concave in $\xi$; it is a \emph{local
information surrogate} and is not claimed to optimize reconstruction PSNR or
any image-domain risk. The estimable subspace is frozen at $r = 200$: $B$ is
the top-$200$ eigenbasis of the fixed comparator's information matrix under
the same pre-scan, dwell, and resource convention, computed before the
optimization and held fixed, with a fixed numerical ridge
$\epsilon_0 = 10^{-9}\,\mathrm{tr}(B^{\top}V_{\text{FIXED}^\ast}B)/r$
(R13~\S6). If every atom obeys a pointwise equal-load constraint the scalar
kernel factors out and the geometry reduces to ordinary linear
D-optimality~\cite{kiefer1959optimum}; genuine dead-time-aware design
therefore requires controlled load variation.

\subsection{The equivalence certificate}
By concavity of $\log\det$ and affineness of $\xi\mapsto V(\xi)$, a feasible
design $\xi^\ast$ is optimal iff there exist constraint multipliers such that
the reduced sensitivity
$s(a) = \mathrm{tr}[V(\xi^\ast)^{-1}H(a)] - \sum_j\alpha_j e_j(a)
- \sum_k\beta_k c_k(a)$ satisfies $s(a) \le \eta^\ast$ over $\mathcal{A}$ with
equality on the support (the constrained Kiefer--Wolfowitz / general
equivalence theorem for nonlinear
models~\cite{kiefer1959optimum,white1973equivalence}). Because a strict
per-pixel-dose dual is not yet implemented, the frozen confirmatory quantity
is the conservative additive budget-dual bound: with
$d_a = \mathrm{tr}(V^{-1} M_{\mathrm{main}} H_a)$, $c_a = \rho_a$, budget
$c^{\top}\xi \le b$, and $\theta \ge 0$,
\begin{equation}
  \begin{aligned}
  G_{\mathrm{rel}} &= \max_a(d_a - \theta c_a) - \sum_a \xi_a(d_a - \theta c_a)
  + \theta\{b - c^{\top}\xi\},\\
  &\qquad\qquad G_{\mathrm{rel}}/r \le 10^{-3},
  \end{aligned}
  \label{eq:relkw}
\end{equation}
which upper-bounds the gap to the optimum of the relaxed (dose-free)
problem and hence to the dose-constrained optimum (R13~\S6.1). It is reported
as \texttt{RELAXED\_KW\_UPPER\_BOUND}, relative to the frozen dictionary and
local estimate, not as a full per-pixel equivalence certificate. The maximum
is taken by an exact scan of the frozen dictionary; a learned proposer may
warm-start the search but the frozen run finishes with the exact
oracle~\cite{jaggi2013frankwolfe}. Continuous weights are rounded to an exact
$972$-row design by efficient rounding followed by deterministic
exchange~\cite{pukelsheim1992rounding,nikolov2022volume,sagnol2015exact},
admissible only if the exact D-efficiency $\ge 0.95$.

\subsection{Photon-budget versus time-budget conjugacy}
\label{sec:conjugacy}
The design admits two conjugate resource corners of the same dead-time
information map. The gate-carrying \textsc{OED-DT} arm holds the original
incident-signal resource as a hard \emph{inequality},
\begin{equation}
  \int \rho_s(a)\,d\xi(a) \le 0.60,
  \label{eq:budget}
\end{equation}
so a subset of high-value atoms may enter the ridge zone only if compensated
by low-load or omitted atoms (R11~\S4); detected counts are dual-reported but
are never the sole fairness budget, because
$\mathbb{E}[N\mid\rho]\simeq\nu\rho/(1+\rho)$ has vanishing marginal cost at
high load. The no-gate \textsc{RIDGE-FIXED} arm instead fixes the support and
applies one global multiplier per dwell, calibrated on the pre-scan estimate
to mean load $\rhoR(\nu)$, followed by a single global safety clip; it uses a
deliberately larger incident budget and is therefore descriptive
(R11~\S2). The exact production ridge is
$\rhoR(\nu) = \min\arg\max_{\rho\in\mathcal{R}_{\mathrm{adm}}(\nu)}
\Jex(\rho,\nu)$ over the admissible set
$\mathcal{R}_{\mathrm{adm}}(\nu) = \{0<\rho\le 24:\ p_{\mathrm{ceil}}\le 0.01,\
J_{\mathrm{mean}}/\Jex \ge 0.90\}$, using the exact numerical ridge rather
than the asymptotic cube-root form (R11~\S1). Under the incident-photon
budget the relevant scalar efficiency is $\Jex(\rho,\nu)/\rho$, which
decreases across the admitted palette; a photon-budgeted optimizer therefore
rationally spends on lower loads and \emph{avoids} the ridge, while the ridge
belongs to the time-limited, power-available corner. This is a
resource-regime statement, not a universal prohibition on ridge-zone
allocation:\ sufficiently valuable directions may still receive larger loads
in other cells (R13~\S8(f)).

\subsection{Guards as constraints}
\label{sec:guards}
The raw optimum is admissible only after all frozen guards pass; each is a
constraint of Eq.~(\ref{eq:master}) or a post-rounding check, not a tunable
proxy. \emph{Peak (two-part, R13~A2):} a shape guard bounds unit-load support
weights to $\tfrac14 \le w_j \le 4$ with support size $k \ge 16$, and a
physical guard caps $\max_{i,j} a^{\mathrm{physical}}_{ij} \le 1536$ (the peak
of a uniform $k=16$ atom at the absolute load cap $\rho = 24$), so
matched-intensity ridge atoms cannot silently demand more peak power.
\emph{Weight power (R10~G4):} the dictionary contains powers $p\in\{0,1\}$
only; $p=2$ is excluded because dev evidence showed higher proxy contrast can
worsen conditioning and endpoint quality. \emph{Per-pixel dose (R10~G5):}
each pixel's mean exposure stays within $\pm5\%$ of the image mean, enforced
in-optimizer and re-verified after rounding; a necessary row bound follows,
$\rho_r \le 1.05\,M\bar d/\max_j g_{rj}$ (R13~A5). \emph{Spectral floor
(R10~G6):} $V_{\mathrm{OED}} \succeq 0.5\,V_{\text{FIXED}^\ast}$ on $B$, with a
convex-mixture fallback whose coefficient is fixed from information matrices
before any image endpoint is seen. \emph{Kernel and trust (R11~G7--G8,
R13~A1):} every $(\nu,\rho)$ node is certified for exact-PMF mass, score, and
derivative agreement, mean-information efficiency
$J_{\mathrm{mean}}/\Jex \ge 0.90$, finite-window RQL log-load bias
$|\log(\rho^{\dagger}/\rho)| \le 0.05$, ceiling probability
$p_{\mathrm{ceil}}(\rho,\nu)\le 0.05$ per atom and $\le 0.01$ design-weighted.
\emph{A-risk (R10~G7):} $\mathrm{tr}(V_{\mathrm{exact}}^{-1}) \le 1.05\,
\mathrm{tr}(V_{\text{FIXED}^\ast}^{-1})$. Any violation is a declared failure
state (\texttt{DICTIONARY\_RANK\_FAIL}, \texttt{CONTINUOUS\_CERTIFICATE\_FAIL},
\texttt{SPECTRAL\_GUARD\_FAIL}, \texttt{A\_RISK\_GUARD\_FAIL},
\texttt{DOSE\_GUARD\_FAIL}, \texttt{ROUNDING\_FAIL}), whose predeclared
fallback is the fixed comparator; no threshold is retuned after freeze.

% ===========================================================================
\section{Theory: the jitter-capped ridge and its unification}
\label{sec:theory}
% R14-RULED (docs/ROUND63_GPT_ROUND14_RULING_RAW.md, issue "R14 ruling:
% unification architecture"): B = headline theorems (random-hold count
% information: identity R14.B1, rate R14.B2, exact cubic optimum),
% A = framework theorem (master program / generalized KKT), C = conditional
% proposition (brightness-information alignment). Corrections frozen by the
% ruling: the smooth crossover REPLACES the earlier min{} conjecture (that
% form is only the two-limit envelope, with an artificial kink); the leading
% coefficient is 2^{-1/3} = 0.794, NOT the fitted 0.72. Exact-law predictions
% were REGISTERED before the refined sweep completed
% (docs/R14_PREREGISTERED_PREDICTIONS.md, commit 8a627bb).
The theory spine follows the frozen unification ruling:\ the random-hold
count-information results of Section~\ref{sec:jitter} are the headline
theorems, the master design program of Section~\ref{sec:onelagrangian} is the
unifying framework, and the brightness--information alignment of
Section~\ref{sec:alignment} is the conditional proposition explaining
\textsc{RIDGE-FIXED}. The deterministic detector of the
companion~\cite{companion} is the zero-jitter critical manifold of the
random-hold theory. In one sentence:\ the recovery-time law determines a
scalar count-information kernel; constrained design allocates illumination
directions and loads against that kernel; and fixed global flux is the
special case in which scene brightness already satisfies the kernel's
water-filling condition. Theorem statements are given here; proofs at
manuscript rigor are deferred to the supplement~\SPH{supplement: proofs of
Theorems~\ref{thm:identity}--\ref{thm:cubic} and
Proposition~\ref{prop:alignment}}.

\subsection{The jitter-capped ridge}
\label{sec:jitter}
Generalize the detector so that after each registered event it blocks for a
hidden \emph{random} hold time $B_i \sim H$, iid, independent of the Poisson
arrivals and of $\lambda$, with mean $\tau$ and coefficient of variation
$c_v = \sigma_B/\tau$ (the companion's deterministic model is $c_v = 0$). The
observed datum remains the integrated registered count $N_T$ over
$T = \nu\tau$; the active time $L_T$, the state path, and the realized holds
are hidden from the count-only observer.

\begin{theorem}[missing-information identity]
\label{thm:identity}
For $\theta = \log\lambda$, the complete detector-path information is
$\mathbb{E}[N_T]$, and the information retained by the scalar count is
exactly
\begin{equation}
  I_N(\theta;T) = \mathbb{E}[N_T]
  - \lambda^{2}\,\mathbb{E}\bigl[\operatorname{Var}(L_T \mid N_T)\bigr]
  \label{eq:identity}
\end{equation}
for every finite $T$ and every iid hidden hold law $H$ independent of
$\lambda$.
\end{theorem}

The companion's terminal-phase identity is the special case in which the
uncertainty in $L_T$ is confined to the terminal detector phase. The identity
requires amendment if the hold law depends on $\lambda$ or the scene, if
realized holds are observed and used, if holds are serially correlated, or if
the counter is paralyzable.

\begin{theorem}[long-window rate]
\label{thm:rate}
Under standard nonlattice renewal local-limit conditions,
\begin{equation}
  \frac{I_N(\log\lambda;T)}{\nu} \;\longrightarrow\;
  J_\infty(\rho,c_v) = \frac{\rho}{(1+\rho)\,(1+c_v^{2}\rho^{2})},
  \label{eq:Jinf}
\end{equation}
depending on $H$ only through its mean and variance. The missing-information
rate converges to $c_v^{2}\rho^{3}/[(1+\rho)(1+c_v^{2}\rho^{2})]$:\ an
\emph{extensive} loss equal, to first order in $c_v^{2}\rho^{2}$, to the
heuristic $\lambda^{2}\sigma_B^{2}\,\mathbb{E}[N_T]$ --- conditioning on the
count partially reveals the accumulated hold fluctuation, contributing the
factor $(1+c_v^{2}\rho^{2})^{-1}$ --- and, unlike the deterministic
terminal-phase loss, nonvanishing as $\nu\to\infty$.
\end{theorem}

\begin{theorem}[exact jitter-limited optimum]
\label{thm:cubic}
For every fixed $c_v>0$, $J_\infty(\rho,c_v)$ has a unique interior maximum
$\rho_c$ solving the cubic
\begin{equation}
  c_v^{2}\,\rho_c^{2}\,(1 + 2\rho_c) = 1,
  \label{eq:cubic}
\end{equation}
with small-jitter expansion
\begin{equation}
  \rho_c = 2^{-1/3}c_v^{-2/3} - \tfrac16 + \tfrac{2^{1/3}}{36}\,c_v^{2/3}
  + O(c_v^{4/3}),
  \label{eq:smallc}
\end{equation}
whose leading coefficient is $2^{-1/3} = 0.79370\ldots$
\end{theorem}

For exponential holds ($c_v = 1$, outside the small-jitter regime) the exact
positive root is $\rho_c = 0.6573\ldots$; gamma, bounded two-point, and
lognormal holds sharing a small $c_v$ share the leading coefficient, with
distribution shape entering at the next order.

\paragraph{Matched crossover (prediction).}
Matching the extensive loss $c_v^{2}\rho^{2}$ against the deterministic
terminal-boundary loss $\rho^{2}/(12\nu)$ in the joint regime $\nu\to\infty$,
$c_v\to0$, $c_v^{2}\nu = O(1)$ gives
\begin{equation}
  \rho_{*}(\nu,c_v) \sim \Bigl(2c_v^{2} + \frac{1}{6\nu}\Bigr)^{-1/3}
  = (6\nu)^{1/3}\,\bigl(1 + 12\nu c_v^{2}\bigr)^{-1/3},
  \label{eq:crossover}
\end{equation}
with scaling variable $x = 12\nu c_v^{2}$ and crossover at
$c_v \asymp (12\nu)^{-1/2}$:\ the deterministic limit $x\to0$ recovers
$(6\nu)^{1/3}$ and the jitter-dominated limit recovers $2^{-1/3}c_v^{-2/3}$.
An earlier $\min\{(6\nu)^{1/3},\, c_0\,c_v^{-2/3}\}$ conjecture is only the
two-limit envelope of Eq.~(\ref{eq:crossover}), with an artificial kink, and
is superseded. Per the frozen claim discipline, the exact identity and the
long-window limit \emph{imply} the jitter exponent $-2/3$; the full
finite-window crossover of Eq.~(\ref{eq:crossover}) is a matched-asymptotic
\emph{prediction} until a uniform two-parameter (triangular-array renewal
local-limit) proof is complete~\SPH{supplement: uniform crossover proof, or
the prediction label is retained}.

\paragraph{Universality class.}
The exponent pair $(1/3, -2/3)$ is universal within the regular
finite-variance iid hidden-hold class:\ iid hidden nonnegative holds whose
law is independent of $\lambda$, with finite nonzero mean and finite
variance, forming a regular small-jitter family, under count-only observation
and active-start nonparalyzable operation. It is \emph{not} claimed for
heavy-tailed (infinite-variance) holds, correlated or drifting holds, holds
observed by the electronics, paralyzable detectors, event-time observations,
or $\lambda$-dependent recovery. The precise critical statement is:\ zero
hold-time variance is the critical manifold on which the optimal load
continues to grow with window length; any fixed nonzero hidden hold variance
produces a finite long-window optimum. Adjacent renewal-CLT and
random-dead-time literature establishes the mean/variance ingredients and
task-specific finite-rate
optima~\cite{muller1973deadtime,yu2000mean,gronberg2018deadtime,
alvarez2014pileup,jorgensen2026deadtime}; we found no source stating the
count-only exponent pair or the extensive/terminal crossover.

% ---- Measured coarse sweep vs. registered exact-law predictions ------------
\begin{table}[t]
\centering
\caption{Monte-Carlo scalar count information $J(\rho,\nu=2000,c_v)$ for
lognormal hold times (coarse sweep:\ common random numbers, $60{,}000$
frames), validated against the exact table at $c_v = 0$ ($J(2) = 0.676$ vs.\
exact $0.667 + \text{corr}$), against the exact-law peak predictions of
Theorem~\ref{thm:cubic}/Eq.~(\ref{eq:crossover}) as registered before the
refined sweep completed (\texttt{docs/R14\_PREREGISTERED\_PREDICTIONS.md},
commit \texttt{8a627bb}). The coarse-grid fit $c_0 \approx 0.72$ is an
effective finite-grid, finite-sample coefficient; the exact asymptotic
coefficient is $2^{-1/3} = 0.794$, and the exact-law predictions sit
$5$--$6\%$ above the coarse measurements (grid-limited; the ratio test
$\mathrm{cap}(0.05)/\mathrm{cap}(0.1) = 1.65$ vs.\ $2^{2/3} = 1.587$ is less
sensitive to the common bias). Hardware warning:\ at $\rho = 24$ with
$c_v = 0.05$, $J = 0.38$ vs.\ deterministic $0.95$ --- about $60\%$
information loss from $5\%$ jitter at the deterministic ridge.
Source:\ \texttt{results/round63\_study2/jitter\_scalar\_fi.log}.}
\label{tab:jitter}
\small
\setlength{\tabcolsep}{4pt}
\begin{tabular}{@{}lccc@{}}
\hline
$c_v$ & measured peak $\rho$ & exact-law prediction & $J$ at peak \\
\hline
$0$    & $\approx 22$--$24$   & $22.25$ (exact ridge) & $0.948$ \\
$0.05$ & $\approx 5.43$       & $5.69$                & $0.796$ \\
$0.1$  & $\approx 3.30$       & $3.50$                & $0.701$ \\
$0.3$  & $\le 2$ (grid edge)  & $1.63$                & $0.499$ \\
\hline
\end{tabular}
\end{table}

The coarse sweep of Table~\ref{tab:jitter} was measured \emph{before} the
exact law was derived; the exact-law column was then registered, together
with two out-of-sample rows --- $c_v = 0.02 \to \rho_{*} \approx 10.4$ and
$c_v = 0.2 \to \rho_{*} \approx 2.30$ --- before the refined sweep completed
(\texttt{docs/R14\_PREREGISTERED\_PREDICTIONS.md}, commit \texttt{8a627bb}).
The refined sweep (three hold families --- lognormal, matched-CV gamma,
bounded two-point --- with cross-fitted conditional-score information
estimation, frozen exponent tests for $1/3$ and $-2/3$, a long-window cubic
residual, and a scaling-collapse test of Eq.~(\ref{eq:crossover})) reports in
Section~\ref{sec:results}.

\subsection{One master concave program}
\label{sec:onelagrangian}
% R14 Candidate A: APPROVED as the framework theorem, with the frozen
% discrete-optima caveat (ruling §2.3) kept verbatim in spirit.
The framework theorem generalizes Section~\ref{sec:design}:\ with the
random-hold kernel $J_H(\rho,\nu)$ of Theorem~\ref{thm:rate} inside the
rank-one atoms $H_z = \nu J_H(\rho_z,\nu)\,q_z q_z^{\top}$, the
design-measure problem of Eq.~(\ref{eq:master}) remains concave, and every
true interior optimum of the declared problem satisfies one generalized
reduced-sensitivity/KKT system:\ the sensitivity
$\psi(z) = M\,\mathrm{tr}[V^{-1}H_z] - \sum_m \beta_m c_m(z)$ obeys
$\psi(z)\le\eta$ with equality on support atoms, and at any support atom
interior in a smooth scalar coordinate $y$, $\partial_y\psi = 0$ with
$\partial^2_y\psi \le 0$. The load coordinate yields the
information-water-filling equation --- $M\nu J_H'(\rho)\,\ell$ equal to the
marginal resource shadow price, with $\ell = q^{\top}V^{-1}q$ --- while the
support-scale and weight-power coordinates equate
$2M\nu J_H(\rho)\,q_y^{\top}V^{-1}q$ to the corresponding marginal
dose/peak/dispersion price:\ the apparent ``conditioning penalty'' is not an
ad hoc term but the $V^{-1}$ geometry inside the reduced sensitivity, which
also explains why raising proxy contrast from $p=1$ to $p=2$ can lower the
true objective. One caveat is kept deliberately:\ the master design problem
places load, support scale, and intensity weighting in one constrained
concave measure optimization, and its support atoms satisfy a common
reduced-sensitivity condition; but the empirically observed occupancy
($k\approx32$) and weighting ($p\approx1$) optima \emph{motivate} these
coordinates and are not themselves claimed as continuous KKT theorems ---
they were selected from discrete grids and judged by image endpoints, not by
the master $\log\det$ objective.

\subsection{Brightness--information alignment}
\label{sec:alignment}
% R14 Candidate C: APPROVED as a conditional proposition, NOT a universal
% approximation theorem. Frozen wording: "without alignment, fixed flux can
% be arbitrarily bad" (no universal constant factor).
At a fixed design iterate with directional values
$\ell_i = q_i^{\top}V^{-1}q_i$, allocate scalar loads under the budget
$\sum_i \pi_i\rho_i = B$ to maximize $F(\rho) = \sum_i \pi_i\ell_i
J(\rho_i)$; on an interval where $J$ is increasing and strictly concave, the
water-filling optimum satisfies $\ell_i J'(\rho_i^{*}) = \beta$ at every
interior coordinate. A single global source induces
$\tilde\rho_i = \kappa u_i$, with $u_i = a_i^{\top}x$ the bucket brightness.

\begin{proposition}[alignment and residual bound]
\label{prop:alignment}
The fixed-flux allocation is exactly optimal for this scalar subproblem if
and only if there exists $\beta$ with
\begin{equation}
  \ell_i\,J'(\kappa u_i) = \beta
  \label{eq:alignment}
\end{equation}
for all interior active directions, with KKT-compatible inequalities at the
bounds. If moreover $-\ell_i J''(\rho) \ge m > 0$ throughout the admitted
interval, then, with residuals $r_i = \ell_i J'(\kappa u_i) - \beta$,
\begin{equation}
  0 \;\le\; F(\rho^{*}) - F(\tilde\rho)
  \;\le\; \frac{1}{2m}\sum_i \pi_i\, r_i^{2}.
  \label{eq:residual}
\end{equation}
\end{proposition}

A computable projected-KKT residual --- not $\Cu$ alone --- therefore
certifies whether one global flux knob is close to the scalar water-filling
optimum. In power-law form, if $J'(\rho)\asymp C\rho^{-s}$ and
$\ell_i \asymp u_i^{\alpha}$, fixed flux aligns when $\alpha\approx s$; for
the rising kernel $J = \rho/(1+\rho)$ the local slope
$s(\rho) = 2\rho/(1+\rho)$ runs from $0$ toward $2$, so the often-invoked
$\ell \propto u^{2}$ alignment is only a high-load approximation. Without the
alignment condition there is no scene-independent constant-factor
guarantee:\ two directions with $u_1/u_2\to\infty$ but $\ell_1/\ell_2\to0$
drive the fixed-flux objective ratio to zero, so fixed flux can be
arbitrarily bad. With random holds, $J'$ vanishes at the jitter cap and is
negative beyond it, so fixed flux can overdrive the brightest directions and
explicit gain control becomes necessary. This explains both sides of the
companion's dev evidence~\cite{companion}:\ \textsc{RIDGE-FIXED} with one
global knob is near-optimal exactly when its realized residuals are small
(self-water-filling), while the uniform servo --- which loses by up to
$\sim7$~dB on sparse-support scenes --- fails because forcing equal $\rho_i$
discards a naturally favorable alignment between $u_i$ and $\ell_i$.

% ===========================================================================
\section{Preregistered campaign (M1)}
\label{sec:campaign}
% Summarized from docs/ROUND63_METHOD_SPEC_M1.md + the R13 freeze audit.
% All outcomes \SPH.
\textsc{M1} is a separate follow-up campaign with its own seed namespace,
manifests, and expected-cell ledger; it does not reopen the companion's
frozen studies. It freezes on the R13 audit plus one immutable commit
(tag \texttt{m1-freeze}); after that commit no scene, atom family, load grid,
guard, endpoint, or gate may change, and if the primary fails there is no
third redesign.

\paragraph{Geometry and constants.}
$n = 32^{2} = 1024$; $M_{\mathrm{total}} = 1024 = M_{\mathrm{pre}}(52) +
M_{\mathrm{main}}(972)$ for every arm; nonparalyzable active start,
$\tau = 50$~ns, $d = 0$, $\sigma_b = 0$;
$\rhobar \in \{0.05, 0.60\}$ for the budget-anchored arms;
$\nu \in \{5, 10, 20, 50, 100, 200, 500, 1000, 2000\}$; five measurement
seeds per (image, condition). Optical time $T_{\mathrm{opt}} =
M_{\mathrm{total}}\,\nu\,\tau$; design-computation wall time is reported
separately and never folded into $T_{\mathrm{opt}}$.

\paragraph{Scenes.}
Twenty-four fresh confirmatory instances (six generator families
$\times$ four instances) with seeds $s_{f,r} = 633000 + 100f + r$, plus six
development-only instances $632900 + f$. The six committed generators (glyph,
chirp, spokes, maze, contour, microtexture) are reused verbatim with their
actual names --- ``USAF'' is not used unless a USAF target is generated, and
spokes is described as a radial resolution target. No acceptance filter of any
kind is applied; generators, images, and ROIs receive SHA256 manifests.

\paragraph{Common pre-scan.}
Every arm receives and pays for the same $52$-pattern scene-independent
balanced pre-scan at its own $(\rhobar,\nu)$:\ $32$ paired $4\times4$-block
rows (support $32$, amplitude $32$), $16$ block rows (amplitude $16$), and
$4$ quadrant rows (amplitude $4$), so every pixel receives cumulative
amplitude $32 + 16 + 4 = 52$ and the row-mean pattern is exactly the all-ones
vector; the mean pre-scan load then equals the declared global load for any
sum-normalized scene, with no $\hat x$-dependent row scaling (R13~\S2). All
$52$ counts enter every arm's final reconstruction; only adaptive arms use the
pre-scan to choose subsequent patterns. Fixed arms run $972$ main patterns;
every arm totals exactly $1024$ physical exposures.

\paragraph{Arms.}
Table~\ref{tab:arms} lists the six arms. Only \textsc{OED-DT} carries a
confirmatory gate; the others are ablations, baselines, or the descriptive
ridge-tracking demonstration. The R13 freeze audit additionally mandates
\textsc{MATCH1} as a no-gate arm and a committed selection of the global fixed
comparator $\text{FIXED}^\ast$ from $\{\textsc{SCAT32}, \textsc{LBLOB16},
\textsc{MATCH1}\}$ by the frozen development rule (median fixed-dwell
radiometric PSNR at $\rhobar=0.60$, $\nu=2000$); \textsc{LBLOB16} is the
current spectral-floor reference~\SPH{M1: committed
$\text{FIXED}^\ast$ selection record}. Reconstruction (RQL with the frozen
analytic $\lambda_{\mathrm{TV}}$ rule and pooled initialization) is identical
across arms; the primary comparison is each arm's own fast curve against the
common safe ($\rhobar = 0.05$) reference of its own pattern set.

\begin{table}[t]
\centering
\caption{The six M1 arms. Only \textsc{OED-DT} carries the confirmatory gate;
the R13 freeze audit additionally requires \textsc{MATCH1} and a committed
$\text{FIXED}^\ast$ selection (see text).}
\label{tab:arms}
\small
\setlength{\tabcolsep}{4pt}
\begin{tabular}{@{}lp{4.0cm}l@{}}
\hline
Arm & Role & Gate \\
\hline
\textsc{OED-DT} & variable-load exact-kernel D-optimal design (v3 solver;
  relaxed-KW certificate vs.\ the frozen dictionary) & primary \\
\textsc{OED-EQLOAD} & same optimizer, all atoms at $\rhobar$ (kernel
  ablation) & none \\
\textsc{SCAT16} & frozen scattered $k=16$ (companion bridge) & none \\
\textsc{SCAT32} & best ladder rung by fixed-dwell median & none \\
\textsc{LBLOB16} & committed clustered support (strongest fixed baseline;
  spectral-floor reference) & none \\
\textsc{RIDGE-FIXED} & dwell-dependent ridge-tracking policy
  (time-budget demonstration) & none \\
\hline
\end{tabular}
\end{table}

\paragraph{Endpoints and gate.}
The frozen $Q_{90}$ time-to-quality machinery transfers verbatim:\ seed-mean
radiometric PSNR, equal-weight PAVA isotonic curves, a $0.50$~dB
safe-range informativeness floor, the frozen censoring taxonomy, and a
$10{,}000$-replicate nested family-stratified bootstrap. The single primary
gate applies to \textsc{OED-DT} only (\texttt{METHOD\_SPEED\_PASS}):\ median
$S_{\mathrm{gate}} \ge 3$, bootstrap lower bound $> 1$, and $\ge 18/24$ scenes
with $S_j > 1$. Two secondaries cannot rescue a failed primary:\ a
fixed-dwell safe-versus-fast gain (\texttt{METHOD\_FIXED\_DWELL\_PASS};
median $\ge 1$~dB at the terminal dwell, lower bound $> 0$, $\ge 18/24$
positive) and a design gain of \textsc{OED-DT} over $\text{FIXED}^\ast$ at the
fast terminal dwell (\texttt{METHOD\_DESIGN\_PASS}; median $\ge 0.50$~dB,
family-stratified lower bound $> 0$, $\ge 18/24$ positive).

\paragraph{Disclosure ledger.}
Every cell discloses incident and detected photons (predicted and realized),
dose ratios, source/peak-irradiance ratio, the load percentiles $\rho_5$,
$\rho_{50}$, $\rho_{95}$, $\rho_{\max}$, predicted and realized ceiling
fractions, and mean $\Jex$; designed cells add the continuous and exact
$\log\det$ objectives, the exact D-efficiency, A-risk, spectral floor, and the
relaxed-KW gap. The full nine-dwell kernel/ridge certification table is frozen
before any confirmatory scene is opened.

% ===========================================================================
\section{Results}
\label{sec:results}
% Skeleton: one \SPH subsection per arm/endpoint + a cross-arm table. Numbers
% fill from the M1 campaign after tag m1-freeze.
\subsection{Primary: OED-DT time-to-quality}
\SPH{M1: METHOD\_SPEED\_PASS verdict --- median $S_{\mathrm{gate}}$, bootstrap
lower bound, count of scenes with $S_j>1$, and the frozen PASS/FAIL call.}

\subsection{Design gain: OED-DT vs.\ $\text{FIXED}^\ast$}
\SPH{M1: METHOD\_DESIGN\_PASS --- median design gain (dB) at the fast terminal
dwell, family-stratified lower bound, count positive.}

\subsection{Fixed-dwell safe-versus-fast gain}
\SPH{M1: METHOD\_FIXED\_DWELL\_PASS --- median terminal-dwell gain (dB),
count positive, bootstrap lower bound, per-arm ladder.}

\subsection{RIDGE-FIXED operating-point demonstration}
\SPH{M1: fixed-dwell gain per dwell of the ridge-tracking policy over
$\rhobar=0.60$ and safe references, realized loads, and ceiling fractions.}

\subsection{OED-EQLOAD kernel ablation}
\SPH{M1: equal-load geometry vs.\ variable-load geometry --- any design
difference attributable to the renewal kernel, with numerical-noise
disclosure.}

\subsection{Certificate and rounding diagnostics}
\SPH{M1: relaxed-KW gap $G_{\mathrm{rel}}/r$, exact D-efficiency, A-risk
ratio, and spectral-floor margin per designed cell.}

\subsection{Jitter-ridge refined sweep}
\SPH{sweep: refined-sweep verdict --- deterministic $1/3$ and jitter $-2/3$
exponent tests, long-window cubic residual, scaling collapse of
$\rho_{*}/(6\nu)^{1/3}$ against $x = 12\nu c_v^{2}$ per
Eq.~(\ref{eq:crossover}), distribution-shape universality across
lognormal/gamma/two-point holds, and the registered out-of-sample rows
$c_v = 0.02 \to \rho_{*}\approx10.4$ and $c_v = 0.2 \to
\rho_{*}\approx2.30$.}

\begin{table*}[t]
\centering
\caption{Cross-arm summary (descriptive; no new gates beyond the
\textsc{OED-DT} primary). All entries fill from the frozen M1 analyzer after
tag \texttt{m1-freeze}.}
\label{tab:crossarm}
\footnotesize
\setlength{\tabcolsep}{5pt}
\begin{tabular}{@{}lccccccc@{}}
\hline
Arm & median $S_{\mathrm{gate}}$ & design (dB) & fixed-dwell (dB)
& mean $\Jex$ & incident budget & ceiling frac. & verdict \\
\hline
\textsc{OED-DT}      & \SPH{M1} & \SPH{M1} & \SPH{M1} & \SPH{M1} & \SPH{M1} & \SPH{M1} & \SPH{M1} \\
\textsc{OED-EQLOAD}  & \SPH{M1} & \SPH{M1} & \SPH{M1} & \SPH{M1} & \SPH{M1} & \SPH{M1} & \SPH{M1} \\
\textsc{SCAT16}      & \SPH{M1} & \SPH{M1} & \SPH{M1} & \SPH{M1} & \SPH{M1} & \SPH{M1} & \SPH{M1} \\
\textsc{SCAT32}      & \SPH{M1} & \SPH{M1} & \SPH{M1} & \SPH{M1} & \SPH{M1} & \SPH{M1} & \SPH{M1} \\
\textsc{LBLOB16}     & \SPH{M1} & \SPH{M1} & \SPH{M1} & \SPH{M1} & \SPH{M1} & \SPH{M1} & \SPH{M1} \\
\textsc{RIDGE-FIXED} & \SPH{M1} & \SPH{M1} & \SPH{M1} & \SPH{M1} & \SPH{M1} & \SPH{M1} & \SPH{M1} \\
\hline
\end{tabular}
\end{table*}

% ===========================================================================
\section{Discussion}
\label{sec:discussion}
% No separate limitations section (companion house style); honesty inside the
% paragraphs. Photon economics + hardware warning + bench outlook.
The two arms expose conjugate resource regimes of the same dead-time
information map. Under a linear incident-photon budget the per-photon
efficiency $\Jex(\rho,\nu)/\rho$ falls monotonically, so a budget-constrained
design rations photons onto lower loads and \emph{avoids} the ridge; the
ridge belongs to time-constrained operation, where power is available and the
only scarcity is dwell. Reading the flux ridge and the design allocation as
one program clarifies that the large fixed-dwell gains of the companion are an
operating-point effect, not evidence that a variable-load optimizer must spend
into the ridge.

The jitter-capped ridge is the practical caution. The cube-root ridge is a
property of the zero-jitter critical manifold; the moment a detector's hold
time carries a coefficient of variation $c_v$, the long-window optimum is
finite --- the exact cubic $c_v^{2}\rho^{2}(1+2\rho) = 1$ of
Theorem~\ref{thm:cubic}, scaling as $2^{-1/3}c_v^{-2/3}$ --- and driving past
it is actively harmful:\ at the deterministic ridge, $5\%$ jitter already
forfeits about $60\%$ of the scalar count information in the coarse dev
measurement (Table~\ref{tab:jitter}), with the final uncertainty of that loss
metric assigned by the preregistered batch protocol of the refined sweep. A
deployment must therefore locate its operating point against the
jitter-limited ceiling of the real detector, not the deterministic ridge; the
safe conclusion is conditional on measured timing jitter.

An optical extension follows the companion's discussion:\ a coarse
complex-field estimate could synthesize an optical displacement that suppresses
the predicted common mode before photon counting, removing photons ahead of
dead-time loss rather than after it~\cite{companion}. All results here are
simulation-based, and every certificate is stated relative to the frozen
dictionary, the local pre-scan estimate, and the declared resource model;
optical-bench validation remains outstanding.

% ===========================================================================
\section{Reproducibility}
\label{sec:repro}
The design formulation, guards, and campaign follow four frozen rulings:\ the
dead-time OED formulation and equivalence certificate (R10), the ridge-zone
load architecture and incident-budget conjugacy (R11), the M1 freeze audit
(R13), and the unification ruling (R14, issue ``R14 ruling: unification
architecture''), whose theorem hierarchy Section~\ref{sec:theory} follows.
Exact-law peak predictions for the refined jitter sweep, including the two
out-of-sample rows, were registered before the sweep completed
(\texttt{docs/R14\_PREREGISTERED\_PREDICTIONS.md}, commit \texttt{8a627bb}).
The campaign freezes at tag \texttt{m1-freeze} once the R13
checklist passes outcome-blind self-tests; the companion is frozen at tag
\texttt{paper1-freeze-v1}. Source-of-record documents for every frozen
constant are cited by ruling number in the block comments of the manuscript
source. All code, frozen preregistration documents, per-cell manifests, and
evidence bundles are available at~\SPH{repo URL wording --- user decision};
funding and acknowledgments are~\SPH{funding/acknowledgments --- user
decision}.

% opticajnl / IEEE port will switch the style; unsrt for the working draft.
\bibliographystyle{unsrt}
\bibliography{refs}

\end{document}
